\documentclass[11pt]{article}

\usepackage[a4paper,margin=29mm]{geometry}
\usepackage{amsmath,amssymb,amsthm}
\usepackage[colorlinks=true,linkcolor=blue,citecolor=blue,urlcolor=blue]{hyperref}
\usepackage{microtype}

\newtheorem{theorem}{Theorem}
\newtheorem{lemma}[theorem]{Lemma}
\newtheorem{corollary}[theorem]{Corollary}

\newcommand{\EG}{Erd\H{o}s--Gy\'arf\'as}

\title{Incidence-Kernel Bounds for Minimal\\
Counterexamples to the \EG{} Conjecture}
\author{Edison Yi\thanks{AI-assisted mathematical research campaign.  The
proof was independently derived and adversarially checked by separate
OpenAI Codex agents; all mathematical claims remain the author's
responsibility.}}
\date{3 August 2026}

\begin{document}
\maketitle

\begin{abstract}
The \EG{} conjecture asserts that every finite graph of minimum degree at
least three contains a cycle whose length is a power of two.  Carr recently
proved that at least $4/7$ of the vertices in a minimal counterexample would
have degree three.  Write $d$ for the number of cubic vertices and $a$ for
the number of vertices of degree at least four.  We prove the uniform bound
\[
 d\ge 2a+6.
\]
Thus every minimal counterexample has at least
$\lceil2|V(G)|/3\rceil+2$ cubic vertices.  The
key object is a graph on the high-degree vertices obtained by suppressing
the cubic vertices with two high-degree neighbours.  Four-cycle exclusion
and minimality make this incidence kernel simple, $C_4$-free, and
$2$-degenerate; a sharp elementary edge bound and parity then give the
result.  We further prove, using one sharp extremal theorem and an
independently replayable finite certificate, that every such counterexample
of order $n\ge34$ satisfies
\[
 d\ge 2a+7,
 \qquad
 d\ge\left\lceil\frac{2n+7}{3}\right\rceil.
\]
Thus equality in the first bound can survive only at the seven orders
\[
15,18,21,24,27,30,33.
\]
\end{abstract}

\section{Introduction}

The \EG{} conjecture asks whether every finite graph of minimum degree at
least three contains a simple cycle of length $2^k$ for some $k\ge2$.
Throughout, graphs are finite, simple, and undirected, and all cycles are
simple.

Following Carr~\cite{Carr2026}, call a counterexample \emph{minimal} if it is
chosen with minimum possible order and, subject to that, minimum possible
size.  Carr proved that at least $4/7$ of its vertices are cubic.  His proof
records two useful consequences of minimality: the vertices of degree at
least four are independent, and the cubic vertices dominate the graph.  We
reprove those observations and retain more of the resulting incidence
structure.

\begin{theorem}\label{thm:main}
Let $G$ be a minimal counterexample to the \EG{} conjecture.  Put
\[
 D=\{v\in V(G):d_G(v)=3\},\qquad
 A=\{v\in V(G):d_G(v)\ge4\},
\]
and write $d=|D|$, $a=|A|$.  Then
\[
 d\ge 2a+6.
\]
Consequently, writing $n=|V(G)|$,
\[
 d\ge \left\lceil\frac{2n+6}{3}\right\rceil
   =\left\lceil\frac{2n}{3}\right\rceil+2,
 \qquad
 a\le \left\lfloor\frac{n-6}{3}\right\rfloor.
\]
\end{theorem}

\section{The incidence kernel}

We begin with the minimality facts, included to make the argument
self-contained.

\begin{lemma}\label{lem:local}
Every proper subgraph $H$ of $G$ has $\delta(H)\le2$.  Moreover, $A$ is
independent and every vertex of $G$ has a neighbour in $D$.
\end{lemma}

\begin{proof}
If a proper subgraph $H$ had minimum degree at least three, it would be a
counterexample of smaller order, or of the same order and smaller size.  If
an edge joined two vertices of $A$, deleting it would preserve minimum degree
at least three.  Hence $A$ is independent.  Finally, if a vertex $v$ had no
neighbour in $D$, deleting $v$ would lower each of its neighbours from degree
at least four to degree at least three, again producing a smaller
counterexample.
\end{proof}

Every vertex of $D$ therefore has at most two neighbours in $A$.  For
$i\in\{0,1,2\}$, let
\[
 D_i=\{v\in D:d_A(v)=i\},\qquad d_i=|D_i|.
\]
Define a graph $J$ on $A$ by placing one edge between the two $A$-neighbours
of each vertex of $D_2$.

\begin{lemma}\label{lem:kernel}
The incidence kernel $J$ is simple, $C_4$-free, and $2$-degenerate.
\end{lemma}

\begin{proof}
The two $A$-neighbours of a vertex in $D_2$ are distinct.  Two different
vertices of $D_2$ cannot have the same neighbour pair, because the two
corresponding length-two paths would form a $C_4$ in $G$.  Thus $J$ is
simple.

Every simple cycle of length $\ell$ in $J$ lifts, by replacing its edges with
their corresponding distinct two-edge paths, to a simple cycle of length
$2\ell$ in $G$.  Hence $J$, and every subgraph of $J$, has no cycle of
power-of-two length.  If a subgraph of $J$ had minimum degree at least three,
it would be a counterexample on at most $a<n$ vertices.  This contradicts
the order-minimality of $G$.  Therefore $J$ is $2$-degenerate.  In
particular it is $C_4$-free.
\end{proof}

We need the following small extremal fact.

\begin{lemma}\label{lem:edges}
If a simple $C_4$-free, $2$-degenerate graph has $r\ge4$ vertices, then it
has at most $2r-4$ edges.  If $r\ge6$, it has at most $2r-5$ edges.
\end{lemma}

\begin{proof}
Every simple $2$-degenerate graph on $r\ge2$ vertices has at most $2r-3$
edges.  Equality forces every deletion in a degeneracy ordering to remove a
vertex of degree exactly two, eventually leaving a triangle.  Reversing the
next deletion adds a vertex adjacent to two vertices of that triangle, which
creates a $C_4$.  This proves the first bound.

Suppose now that $r\ge6$ and equality $|E|=2r-4$ holds in the first bound.
The same deletion argument reduces an equality example to the unique
four-vertex, four-edge $C_4$-free graph, the paw, and then to the unique
five-vertex, six-edge example, two triangles sharing one vertex.  Every pair
of vertices in the latter graph has a common neighbour.  Adding the next
degree-two vertex therefore creates a $C_4$, a contradiction.
\end{proof}

\section{Preliminary abundance bounds}

\begin{lemma}\label{lem:preliminary}
Every minimal counterexample satisfies $d\ge2a+4$.  If $a\ge6$, then
$d\ge2a+6$.
\end{lemma}

\begin{proof}
First assume $a\ge4$.  Since the edges of $J$ correspond bijectively to the
vertices of $D_2$, Lemma~\ref{lem:edges} gives
\[
 d_2\le2a-4.
\]
Let $e=e_G(A,D)$ and introduce the nonnegative incidence slacks
\[
 x=e-4a,
 \qquad
 y=2d-e=d_1+2d_0.
\]
Since $e=d_1+2d_2$, we have $d_1\ge8+x$ and hence $y\ge8+x$.  The exact
identity
\[
 2(d-2a)=x+y
\]
now yields $d\ge2a+4$.

If $a\ge6$, the second part of Lemma~\ref{lem:edges} instead gives
$d_2\le2a-5$, so $d_1\ge10+x$ and $y\ge10+x$.  Thus $d\ge2a+5$.  Equality
would force $x=0$ and $y=10$.  But the degree sum inside $G[D]$ is
\[
 \sum_{v\in D}d_{G[D]}(v)=3d-e=d+y,
\]
which would be odd when $d=2a+5$ and $y=10$.  This contradicts the
handshake lemma, and proves $d\ge2a+6$.

It remains to establish the universal bound for $a\le3$.
If $a=0$, then $G$ is a nonempty simple cubic graph, so $d\ge4$.  If $a=1$
and $d\le5$, the sole $A$-vertex has at least four distinct neighbours in
$D$.  For $d=4$, the graph $G[D]$ is a four-cycle.  For $d=5$, parity forces
all five vertices of $D$ to meet $A$, so $G[D]$ is a five-cycle; the
$A$-vertex together with any two-edge segment of that cycle gives a
$C_4$.  Hence $d\ge6$.

If $a=2$, simplicity and $C_4$-freeness give $d_2\le1$.  Thus
\[
 8\le d_1+2d_2\le d+1.
\]
The only value below $2a+4=8$ is $d=7$, which forces
$(d_0,d_1,d_2)=(0,6,1)$.  The degree sum in $G[D]$ is then
$2d_1+d_2=13$, impossible.  Similarly, if $a=3$, then $d_2\le3$ and
\[
 12\le d_1+2d_2\le d+3.
\]
The only value below $2a+4=10$ is $d=9$, forcing
$(d_0,d_1,d_2)=(0,6,3)$ and the odd internal degree sum
$2d_1+d_2=15$.  This completes the proof.

\end{proof}

\section{Excluding the remaining equality profiles}

It remains to rule out $d=2a+4$ and $d=2a+5$ when $a\le5$.  Write
$R=G[D]$ and colour each vertex of $D_1$ by its unique neighbour in $A$.

\begin{lemma}\label{lem:twostep}
If two vertices of $D_1$ are the endpoints of a two-edge path in $R$, then
their colours are distinct and have no common neighbour in $J$.
\end{lemma}

\begin{proof}
If $u-z-w$ is such a path and $u,w$ have the same colour $\alpha$, then
$\alpha,u,z,w$ form a four-cycle.  Suppose instead that their distinct
colours $\alpha,\beta$ have a common $J$-neighbour $\gamma$.  Let
$p,q\in D_2$ encode the kernel edges $\alpha\gamma$ and
$\beta\gamma$.  Then
\[
 \alpha,u,z,w,\beta,q,\gamma,p,\alpha
\]
is an eight-cycle.  Its vertices are distinct; in particular $z$ cannot be
$p$ or $q$, because vertices of $D_2$ have degree one in $R$.
\end{proof}

For $s=d-2a\in\{4,5\}$, the incidence identities take the form
\begin{equation}\label{eq:profiles}
 x+y=2s,\qquad
 d_2=2a-s+d_0+x,\qquad
 d_1=2s-x-2d_0.
\end{equation}
The degree sum $d+y$ in $R$ is even.  Thus $x$ is even when $s=4$ and odd
when $s=5$.

\begin{proof}[Completion of the proof of Theorem~\ref{thm:main}]
Lemma~\ref{lem:preliminary} already excludes both values of $s$ for
$a\ge6$.  We treat $a\le5$.

If $a=0$, the case $s=4$ is the cubic graph on four vertices, namely $K_4$,
and the case $s=5$ contradicts the handshake lemma.  Both are impossible.

Let $a=1$.  Since $d_2=0$, equation~\eqref{eq:profiles} leaves, for each
$s$, either a $2$-regular graph $R$ all of whose vertices have the same
colour, or $d_0=2$ and respectively four or five vertices in $D_1$.  The
first alternative contains a two-edge segment closing through the sole
$A$-vertex to a $C_4$.  In the second, the two-element $R$-neighbourhoods of
the $D_1$ vertices must be pairwise disjoint by
Lemma~\ref{lem:twostep}; respectively eight elements cannot fit in six
vertices, and ten cannot fit in seven.

Let $a=2$.  Since $J$ is simple, equation~\eqref{eq:profiles} leaves two
profiles for each $s$.  If $J$ is empty, then $R$ is $2$-regular on eight
or nine coloured vertices.  On a cycle of length $\ell$, adjacency at
distance two gives one odd cycle for odd $\ell$ and two cycles of length
$\ell/2$ for even $\ell$.  It is bipartite only when $4$ divides $\ell$.
Avoiding a same-coloured distance-two pair would therefore make every
component length divisible by four, forcing a forbidden $C_4$ or $C_8$.

In the other profile, $J$ is one edge and $R$ has one degree-three vertex
$z$, one degree-one vertex $e\in D_2$, and all remaining vertices in $D_1$
of degree two.  Its non-cycle component is unicyclic, consisting of a cycle
through $z$ and a pendant path from $z$ to $e$; any other component would be
a coloured cycle and is excluded by the preceding distance-two argument.
Since $e$ is adjacent to both vertices of $A$, the pendant path has exactly
one coloured internal vertex: zero or at least two would immediately give a
$C_4$.  For $s=4$, the cycle is a six-cycle containing five coloured
vertices; three alternating coloured vertices are pairwise at distance two,
so two have the same colour and give a $C_4$.  For $s=5$, write the cycle as
\[
 z,c_1,c_2,c_3,c_4,c_5,c_6,z.
\]
The distance-two constraints on its coloured vertices form the path
$c_2,c_4,c_6,c_1,c_3,c_5$.  If the two colour classes on these six vertices
are unbalanced, this path is not properly coloured and gives a $C_4$.  If
they are balanced, proper colouring forces alternation, so $c_2,c_3$ have
the same colour.  Their complementary six-edge path around the seven-cycle,
together with the two edges through their common $A$-neighbour, is a
$C_8$.

It remains to treat $a=3,4,5$.  Equations~\eqref{eq:profiles}, simplicity,
and the kernel edge bounds leave exactly the kernels $P_3$, $K_3$, the paw,
and the five-vertex friendship graph $F_2$.  For $s=4$ their values of
$d_1$ are respectively $8,6,8,8$; for $s=5$ they are $9,7,9,9$.

For $J=P_3$, call the colours $L,M,R$ in path order.  Let $h$ be the number
of $D_1$ vertices with two $D_1$ neighbours in $R$.  The two degree-one
vertices of $D_2$ give
\[
 h\ge6\quad(s=4),\qquad h\ge7\quad(s=5).
\]
Lemma~\ref{lem:twostep} forces the two neighbour colours of each such vertex
to be $M$ and one of $L,R$.  There are at most two $M$-coloured vertices for
$s=4$, and at most three for $s=5$, with total $R$-degree at most four and
six respectively.  Both inequalities are impossible.

For $J=K_3$, Lemma~\ref{lem:twostep} says no vertex of $R$ has two
$D_1$ neighbours.  Each of the six or seven vertices of $D_1$ therefore
needs a neighbour outside $D_1$.  The three degree-one vertices of $D_2$
supply three incidences, while the sole degree-three vertex of $D_0$ can
supply at most one without becoming the middle of a forbidden two-edge
path.  The total capacity four is insufficient.

For the paw, the four degree-one vertices of $D_2$ imply
$h\ge4$ for $s=4$ and $h\ge5$ for $s=5$.  The only distinct colour pair
without a common $J$-neighbour is the degree-three centre and the leaf.
Thus every counted vertex must be incident with a centre-coloured vertex.
There are respectively at most one and two such vertices, of total
$R$-degree at most two and four, again a contradiction.

Finally, in $F_2$ any two distinct colours occurring on $D_1$ have a common
$J$-neighbour, including the shared centre if it receives the one extra
colour in the $s=5$ profile.  Hence no vertex of $R$ has two $D_1$
neighbours.  All eight or nine degree-two vertices of $D_1$ require a
neighbour in $D_2$, but the six degree-one vertices of $D_2$ supply only six
incidences.

We have excluded $s=4,5$ in every remaining case.  Together with
Lemma~\ref{lem:preliminary}, this proves $d\ge2a+6$.  Substituting
$a=n-d$ gives $3d\ge2n+6$ and $3a\le n-6$, which are exactly the stated
integer bounds.
\end{proof}

\section{Excluding the top equality layer}

We now sharpen Theorem~\ref{thm:main} outside a finite window.  This section
uses the following result of Gy\H{o}ri, Li, Salia, Tompkins, Varga and
Zhu~\cite{GLSTVZ2026}: every $r$-vertex graph with more than
\[
 \left\lfloor\frac{19(r-1)}{12}\right\rfloor
\]
edges contains a cycle whose length is divisible by four.

Suppose temporarily that equality holds in Theorem~\ref{thm:main}, so that
\begin{equation}\label{eq:plus-six-equality}
 d=2a+6.
\end{equation}
Put
\[
 q=2a-|E(J)|=2a-d_2.
\]
Using the incidence slacks $x,y$ from Lemma~\ref{lem:preliminary}, direct
substitution gives
\begin{equation}\label{eq:q-six-profiles}
 d_2=2a-6+x+d_0,\qquad
 q=6-x-d_0,\qquad
 d_1=12-x-2d_0.
\end{equation}
When $a\ge6$, Lemma~\ref{lem:edges} gives $q\ge5$.  Moreover, $d$ is even,
and the degree sum $d+y$ inside $G[D]$ shows that $y$, and hence $x$, is
even.  Therefore exactly two profiles remain:
\begin{align}
 q=5 &: \quad x=0,\ d_0=1,\ d_1=10,\label{eq:q5}\\
 q=6 &: \quad x=0,\ d_0=0,\ d_1=12.\label{eq:q6}
\end{align}
In both profiles every vertex of $A$ has degree four, and the number of
$D_1$ vertices whose unique neighbour in $A$ is $v$ equals
\begin{equation}\label{eq:colour-multiplicity}
 c(v)=4-d_J(v).
\end{equation}

For completeness, the extremal theorem makes both profiles finite.  A
dyadic-cycle-free, 2-degenerate graph of deficit five,
$|E(H)|=2|V(H)|-5$, is connected once its order is at least five.  From
order seven onward it has no degree-one vertex; 2-degeneracy therefore
supplies a degree-two vertex whose deletion preserves the deficit.  Thus an
example of order at least nine would reduce to order nine, where it has
thirteen edges, exceeding
$\lfloor19(9-1)/12\rfloor=12$.  The resulting cycle divisible by four has
length four or eight, a contradiction.  Hence the profile~\eqref{eq:q5}
has $a\le8$.

The same deficit addition shows that a disconnected deficit-six graph has
order at most six.  Once deficit five is excluded from order nine onward, a
deficit-six graph of order at least eleven has no degree-one vertex and
again deletes through degree two.  At order eleven its sixteen edges exceed
$\lfloor19(11-1)/12\rfloor=15$.  Hence the profile~\eqref{eq:q6} has
$a\le10$.  It follows already that equality~\eqref{eq:plus-six-equality} is
impossible for $a\ge11$.

It remains to exclude $a=10$ in the deficit-six profile.  Here $J$ is a
connected ten-vertex, fourteen-edge graph with
\[
 2\le\delta(J)\le\Delta(J)\le4,
\]
which is 2-degenerate and has no $C_4$ or $C_8$.  Connectedness follows by
adding the component deficits; a degree-one vertex would delete to the
already excluded nine-vertex deficit-five case.  Exact canonical generation
leaves precisely the following four graph6 records:
\begin{verbatim}
I?`D@`keO
I?`D@pScg
I?`DCdcMG
I?`@dEWR_
\end{verbatim}
All four have degree sequence
$2,2,2,2,3,3,3,3,4,4$.

We describe the independently replayable finite obstruction.  Let
$R=G[D]$.  Each of the fourteen vertices in $D_2$ has degree one in $R$,
each of the twelve vertices in $D_1$ has degree two, and $D_0$ is empty.
Thus $R$ is a disjoint union of exactly seven paths with $D_2$ endpoints,
together with cycles consisting only of $D_1$ vertices.  If
$\lambda_1,\ldots,\lambda_7$ are the numbers of internal $D_1$ vertices on
the paths and $\mu_1,\ldots,\mu_s$ are the cycle lengths, then
\begin{equation}\label{eq:166-topologies}
 \lambda_i\ge0,\qquad
 \mu_j\in\{3,5,6,7,9,10,11,12\},\qquad
 \sum_i\lambda_i+\sum_j\mu_j=12.
\end{equation}
There are exactly 166 unordered topologies in~\eqref{eq:166-topologies}.

Subdivide every edge of $J$ by its corresponding $D_2$ vertex to obtain the
fixed incidence skeleton.  If a segment of $R$ and a simple path in this
skeleton have the same endpoints, their internal vertices lie respectively
in $D_1$ and in $A\cup D_2$.  Their union is therefore a simple cycle, so
the sum of their lengths cannot belong to $\{4,8,16,32\}$.

The verifier enumerates every simple skeleton path and every locally legal
path or cycle component.  A component is represented by its set of $D_2$
endpoint labels and its color-count vector.  Exact set-packing dynamic
programming then requires the endpoint sets to partition all fourteen
labels and the color counts to equal~\eqref{eq:colour-multiplicity}.  This
search is a relaxation: it imposes only the one-segment necessary
conditions, so an empty search space is sufficient for exclusion.

For the first two kernels, paths with respectively zero, one and two
internal $D_1$ vertices have $91,8,25$ legal signatures; for the last two
the counts are $91,0,20$.  Every longer path table and every cycle table is
empty.  The two nominal length-one paths in the first two kernels reuse the
same endpoint pair, leaving only topology $(0,2,2,2,2,2,2)$.  In that
topology the union of the color supports of all legal length-two paths omits
a color whose multiplicity in~\eqref{eq:colour-multiplicity} is positive,
for each of the four kernels.  Thus no residual assignment exists.

The standard-library verifier
\texttt{audit\_a10\_residual\_gate\_2026-08-03.py} independently regenerates
the four kernels and the 166 topologies, enumerates the component signatures,
and performs the exact packing.  Its output status is
\texttt{VERIFIED\_COMPLETE\_NO\_ASSIGNMENT}; the canonical output and an
independent replay are byte-identical with SHA-256
\begin{center}
\small\texttt{F3E39B881120A2B0FFDD7BCFE76E6778FDD2A38D2D3F66AE59745367E4767A25}.
\end{center}

At the opposite end, equality is also impossible for $a\le2$.  We use the
following elementary finite lemma.

\begin{lemma}\label{lem:hamilton-twelve}
Every Hamiltonian graph on twelve vertices with minimum degree at least
three contains a $C_4$ or a $C_8$.
\end{lemma}

\begin{proof}
Fix a Hamilton cycle and label its vertices alternately
$E_0,O_0,\ldots,E_5,O_5$, with subscripts modulo six.  A chord at cyclic
distance $3,5,7$, or $9$ immediately closes to a $C_4$ or a $C_8$.
Consequently every remaining chord joins vertices of the same parity.  Since
every vertex needs a chord, the chords in each parity class contain an edge
cover of six vertices.

Choose an inclusion-minimal such cover.  It is a star forest, so its
component-order partition is $6$, $4+2$, $3+3$, or $2+2+2$.  Comparing the
two Hamilton-cycle arcs between pairs of retained chords eliminates the
first three cases.  Of the fifteen perfect matchings in the last case,
exactly the following six survive within one parity class:
\[
\begin{array}{lll}
01|23|45,&01|25|34,&03|12|45,\\
03|14|25,&05|12|34,&05|14|23.
\end{array}
\]
Combining one surviving even matching with one surviving odd matching gives
a $C_4$ or a $C_8$ in each of the $36$ cases.  The standard-library script
\texttt{audit\_hamilton12\_lemma.py} independently enumerates the minimal
edge covers and prints a literal cycle witness for all $36$ pairs.
\end{proof}

Under equality, degree summation gives
\[
 n=3a+6,\qquad |E(G)|=5a+9+x/2,
\]
where $x\ge0$ is even.  For $a=0,1,2$, the extremal thresholds are
respectively $7,12,17$, while $G$ has at least $9,14,19$ edges.  The forced
cycle divisible by four has length four at order six, length four or eight
at order nine, and length four, eight, or twelve at order twelve.  Only the
last alternative survives, but it is Hamiltonian and
Lemma~\ref{lem:hamilton-twelve} supplies a $C_4$ or a $C_8$.  Thus equality
is impossible for $a\le2$.

We have proved the following strengthening.

\begin{theorem}\label{thm:plus-seven}
Let $G$ be a minimal counterexample, with $a,d,n$ as above.  If
$d=2a+6$, then $a\le9$.  Consequently, if $n\ge34$, then
\[
 d\ge2a+7,
 \qquad
 d\ge\left\lceil\frac{2n+7}{3}\right\rceil,
 \qquad
 a\le\left\lfloor\frac{n-7}{3}\right\rfloor.
\]
Equality in Theorem~\ref{thm:main} can survive only at the seven orders
$15,18,21,24,27,30,33$.
\end{theorem}

\begin{proof}
The deficit analysis excludes equality for $a\ge11$, the finite certificate
excludes $a=10$, and the small-end argument excludes $a\le2$.  Hence equality
implies $3\le a\le9$.  If $n\ge34$ and equality held, then
$n=a+d=3a+6$ would force $a\ge10$, a contradiction.  Thus $d\ge2a+7$;
substituting $a=n-d$ gives the two displayed integer bounds.
\end{proof}

\section{A related size bound}

The same minimality framework gives a useful accompanying estimate.

\begin{corollary}\label{cor:edges}
If $m=|E(G)|$, then $m\le2n-3$.
\end{corollary}

\begin{proof}
Choose a cubic vertex $v$.  Every subgraph of $G-v$ is a proper subgraph of
$G$, so $G-v$ is $2$-degenerate.  Hence
\[
 m-3=|E(G-v)|\le2(n-1)-3,
\]
which gives $m\le2n-2$.

If equality held, $G$ would be an $n$-vertex, $(2n-2)$-edge graph with no
proper subgraph of minimum degree three.  Narins, Pokrovskiy and
Szab\'o~\cite[Theorem~1.4]{NPS2017} proved that every such graph is
pancyclic.  In particular it contains a $C_4$, a contradiction.
\end{proof}

\section{Verification and priority boundary}

The proof was independently derived and then checked line by line by a
separate agent.  The kernel edge bounds were also exhaustively tested on
every labelled graph through seven vertices; this computation is an audit
aid, not a premise.  The finite certificate for
Theorem~\ref{thm:plus-seven} was replayed independently from the exploratory
solver, and its reduction and implementation received a separate hostile
audit.  Searches on 3 August 2026 for the exact constants and the
incidence-kernel construction found Carr's $4/7$ theorem but no earlier
statement of the bounds above.  That is a search-relative novelty check, not
a guarantee concerning unpublished or unindexed work.  These theorems are
necessary conditions on a hypothetical counterexample and do not resolve
the \EG{} conjecture.

\paragraph{Acknowledgement.}
We thank Avery Carr for the structural observations that motivated this
refinement.  The result was developed in an AI-assisted research campaign;
all statements were retained only after adversarial quantifier-level review.

\enlargethispage{4\baselineskip}
\begin{thebibliography}{9}
\small

\bibitem{Carr2026}
A. Carr,
\emph{Every Minimal Counterexample to the Erd\H{o}s--Gy\'arf\'as Conjecture
is Predominantly Cubic},
arXiv:2605.22844, 2026.
\url{https://arxiv.org/abs/2605.22844}.

\bibitem{NPS2017}
L. Narins, A. Pokrovskiy and T. Szab\'o,
\emph{Graphs without proper subgraphs of minimum degree 3 and short cycles},
Combinatorica \textbf{37} (2017), 495--519.
\url{https://arxiv.org/abs/1408.5289}.

\bibitem{GLSTVZ2026}
E. Gy\H{o}ri, B. Li, N. Salia, C. Tompkins, K. Varga and M. Zhu,
\emph{On graphs without cycles of length $0$ modulo $4$},
Journal of Combinatorial Theory, Series B \textbf{176} (2026), 7--29.
\url{https://arxiv.org/abs/2312.09999}.

\end{thebibliography}

\end{document}
